\subsection{Candidate pools and exact enumeration}\label{sec:p2pools}
La and Ti, in that site order, minimize predicted violation in both candidate pools. We fix $X=\mathrm{NO_2}$, the dataset label for two O atoms and one N atom. For each ordered cation pair $(a,b)$, its cost $W_{ab}$ is the mean prediction of the 20 H-soft members fitted using files 02 and 03. The same predictors supply every cost, including costs for compositions with known training labels.

Figure~\ref{fig:p2pipeline} shows which pairs remain eligible. File 07 contains all $52\times52$ pairs. File 11 contains the 529 pairs in this slice of the unlabeled test file 04; the other 2,175 pairs occur in files 02 or 03. The blind mask has holes, so selecting an A and a B independently can produce a forbidden pair.
\paperfigure{P2/plates/P2-F1_data_pipeline.pdf}{Predicted violation over the two candidate pools. Hatching marks excluded pairs. Steps give observed bin shares within each pool; the curves interpolate the bin centres. The star marks LaTiO$_2$N.}{fig:p2pipeline}

Table~\ref{tab:p2pools} gives the nearest competitor in each pool. Here and below, original regret is a candidate's predicted violation minus the least predicted violation among allowed candidates in the same pool.
\begin{table}[!htb]
\centering
\caption{Exact minima and nearest competitors on the frozen surrogate table}\label{tab:p2pools}
\begin{tabular}{@{}lrlrlr@{}}
\toprule
Pool & Pairs & Minimum & $\widehat V_{\min}$ & Runner-up & Gap\\
\midrule
Full (07) & 2,704 & LaTiO$_2$N & 1.839047 & ScTiO$_2$N & 0.416697\\
Blind (11) & 529 & LaTiO$_2$N & 1.839047 & SrTaO$_2$N & 0.465139\\
\bottomrule
\end{tabular}
\end{table}

LaTiO$_2$N has an experimental counterpart. Kasahara et al. measured hydrogen and oxygen evolution in separate visible-light reactions with a sacrificial donor and acceptor, respectively~\cite{kasahara2002}. The blind pool supplies no joint measurement of the five contest criteria for this candidate. Figure~\ref{fig:p2crystal} shows the Ti coordination in an external structure~\cite{mp1222738}.
\paperfigure{P2/plates/P2-A1_crystal_orientations.pdf}{Three views of the same LaTiO$_2$N supercell, with matched ball-and-stick views below. Materials Project mp-1222738~\cite{mp1222738}, accessed 28 November 2026. Illustrative VESTA rendering, not DFT-relaxed in this study. All views derive from one structure and supply no independent observations. Its coordinates served only for display and did not enter prediction, encoding or selection.}{fig:p2crystal}
\FloatBarrier

\subsection{QUBO formulation and integer encoding}\label{sec:p2encoding}
Let $m$ and $m+R$ be the least and greatest allowed entries of $W$ in the pool. Set $C_{ab}=W_{ab}-m$ on allowed pairs and $C_{ab}=R$ elsewhere. The two binary registers $x$ and $y$ select the A and B elements. The quadratic unconstrained binary optimization (QUBO) objective is
\begin{equation}
 H(x,y)=\sum_{a,b} C_{ab}x_a y_b+
 \lambda\left[\left(\sum_a x_a-1\right)^2+
                 \left(\sum_b y_b-1\right)^2\right],\qquad \lambda=1.1R .
 \label{eq:p2qubo}
\end{equation}\nopagebreak[4]
\noindent Each squared term vanishes when its register selects exactly one element. All costs are nonnegative. An allowed minimum attains zero, a forbidden one-hot choice costs $R>0$, and a register with the wrong count incurs at least $\lambda$. Thus the floating objective selects precisely the allowed minimum.

For the row--column representation, subtract each row minimum $u_a$ from $C_{ab}$, then each residual column minimum $v_b$. Put $u_a x_a$ and $v_b y_b$ back as linear terms. Every one-hot pair keeps its original lookup cost. For both representations, binary variables map to spins via $x=(1+s s_0)/2$, with a shared reference spin $s_0$. We use $E(s)=-s^{\mathsf T}Js$ for a symmetric matrix $J$ with zero diagonal. The one-hot comparison scales the largest absolute entry to 127 and rounds every coefficient to the nearest integer. Both use the same penalty. The strongest blind-pool coefficient joins a register spin to the reference (absolute weight 964.56 direct, 902.37 row--column). Cross-register weights reach only 6.01 in magnitude; after rounding, each becomes either zero or $-1$.

\begin{proposition}\label{prop:p2onehot}
Neither audited 8-bit one-hot matrix guarantees that every integer ground state is an allowed blind-pool pair.
\end{proposition}
\begin{p1proof}
Enumerating the 529 allowed pairs gives minimum energies $-10822$ and $-10994$ for the direct and row--column matrices, respectively. In the direct matrix, a forbidden one-hot pair attains $-10822$. In the row--column matrix, a vector with the wrong register counts attains $-10994$. In each case an illegal vector is at or below every legal vector. A still lower ground state would itself have to be illegal.
\end{p1proof}
Among the allowed integer minimizers in Proposition~\ref{prop:p2onehot}, the worst original regrets are 21.65085 and 13.55685, respectively. Rounding can therefore affect selection even when decoding is restricted to allowed pairs. Appendix Figures~\ref{fig:p2graph} and~\ref{fig:p2precision} give the register couplings and precision controls.

\textbf{Encoding the allowed dictionary.} Index the 529 allowed cells by row, then column, starting at zero. Let $r_i=W_i-m$ be candidate $i$'s original regret. The integer cost $q_i$ gives the low-violation end of the table three quarters of the available scale:
\begin{equation}
 q_i=\operatorname{round}(63 f(r_i)),\qquad
 f(r)=\begin{cases} 3r/8, &0\le r\le2,\\
 3/4+(r-2)/(4(R-2)), &2<r\le R .
 \end{cases}\label{eq:p2shape}
\end{equation}\nopagebreak[4]
\noindent Here $R=48.07978$ for the blind pool. The continuous map increases strictly. Rounding can create ties, so we check all integer costs directly: only LaTiO$_2$N has $q_i=0$. Construction uses the full table, including its known minimum.

Divide the row-ordered dictionary into 40 contiguous groups of 13 or 14 candidates. A group of $n+1$ candidates uses $n$ binary variables $z_j$. The legal word with $t$ leading ones followed by zeros selects candidate $t$ in that group, counting from zero. This domain-wall representation~\cite{chancellor2019} has the following group QUBO:
\begin{equation}
 H_g=q_{g,0}+\sum_{j=1}^{n}(q_{g,j}-q_{g,j-1})z_j
       +128\sum_{j=1}^{n-1}(1-z_j)z_{j+1}.
 \label{eq:p2chain}
\end{equation}\nopagebreak[4]
\noindent The successive differences telescope to $q_{g,t}$ on a legal word. Each forbidden transition from zero to one adds a penalty of 128. Packing these 40 distinct subproblems uses $529-40=489$ chain spins and one shared reference. The largest absolute integer coupling is 73. Each legal group supplies one candidate; we select the smallest pair $(q_i,i)$ across groups, using at most 39 classical comparisons. Figure~\ref{fig:p2coupling} shows how the two encodings distribute their couplings.
\paperfigure{P2/plates/P2-F2_coupling_matrices.pdf}{Blind-pool couplings. The left panel shows the floating logical one-hot matrix, the right the submitted integer matrix. Both include the reference spin. Blue and amber denote negative and positive entries on the indicated symmetric logarithmic scales. Every nonzero entry is shown (threshold zero). The layouts describe model couplings and carry no physical SPQC connectivity information.}{fig:p2coupling}
\FloatBarrier

\subsection{Exact, annealing and cloud solvers}\label{sec:p2solve}
\textbf{Enumeration.} Scanning the allowed cells gives Table~\ref{tab:p2pools} directly. It also provides the original-unit reference for every decoded solution. The cloud evaluation covers the blind pool alone.

\textbf{Local simulated annealing (SA).} We use \texttt{Simulated\allowbreak Annealing\allowbreak Optimizer} from the classical module of Kaiwu SDK 1.4.4. It cools from 26,284 to 2.6284 by a factor of 0.99 per step, with 550 iterations per temperature. Each seeded restart retains one solution. Five fixed blocks of 100 restarts provide the local comparison.

\textbf{Cloud optimization.} The QBoson Kaiwu client \texttt{cim.CIMOptimizer} submitted the same matrix through its default route in \texttt{optimization} mode~\cite{kaiwu144}. The service controls its internal effort, ranks and clusters solutions, and returned ten selected vectors per task. Its physical draw count is unavailable. P01--P05 denote five pilot tasks used to check the service; F01--F20 denote the twenty formal tasks. All 25 have distinct result IDs. The matrix and formal selection rule were fixed before F01.

For each vector we recompute $E$ from the returned spins and require all 40 chains to be legal. A task attains the proxy optimum if at least one strict vector decodes to zero original regret. A separate rule selects just one strict vector by lowest integer energy, breaking ties by return order. Scoring that one vector on $W$ tests whether energy alone suffices for the choice. Appendix Figures~\ref{fig:p2attain} and~\ref{fig:p2time} document the endpoints and observation windows.

\subsection{Certificates and residual bounds}\label{sec:p2cert}
\begin{proposition}\label{prop:p2ground}
For the submitted 490-spin matrix, every integer ground state is legal and decodes to LaTiO$_2$N. The ground energy is $-58722$ and the worst ground-state original regret is zero.
\end{proposition}
\begin{p1proof}
Fix the common reference spin to $+1$ by global sign symmetry. Each group can then be minimized independently. In each group a two-state spin recursion retains the least prefix energy for each last spin and a flag for any forbidden transition. Adding spin $t$ after $s$ at position $j$ adds $-2t(J_{j,\mathrm{ref}}+J_{j-1,j}s)$, where the reference is shared by all groups; the first spin contributes only its reference term. This exhausts the paths in each group. Summing the legal minima gives $-58722$. The least cost of replacing any one legal group minimum with an illegal one is 340. Thus every ground state is legal. On legal words the integer energy of each group is a group constant plus $4q_{g,t}$. The groups cover the dictionary exactly once, and group 17 alone contains the unique zero-cost candidate. At a ground state its selected candidate therefore wins the classical comparison.
\end{p1proof}

\begin{proposition}\label{prop:p2margin}
Let $\Delta=E-E_0$ denote the integer energy above the certified ground energy $E_0$. Every vector with $\Delta<196$ is fully legal and selects the unique proxy optimum. Equality admits a wrong choice.
\end{proposition}
\begin{p1proof}
Proposition~\ref{prop:p2ground} gives an illegal-state gap of 340, ruling out a broken chain below 196. For a legal vector with group choices $t_g$, the energy difference is
\begin{equation}
 \Delta=4\sum_{g=1}^{40}\big(q_{g,t_g}-\min_t q_{g,t}\big).
 \label{eq:p2excess}
\end{equation}\nopagebreak[4]
\noindent Each summand is nonnegative. In group 17 the next lowest cost after zero is 49, so losing LaTiO$_2$N costs at least $4\times49=196$. At equality one can choose dictionary entry 217 in that group and local minima in the other 39. The classical comparison then selects entry 448, SrTaO$_2$N, with original regret 0.465139. This attains the boundary.
\end{p1proof}

Figure~\ref{fig:p2margin} applies this matrix bound to the observed returns. F03 had the least headroom among the 20 energy-selected solutions (20 units). The smallest margin among the 200 formal returned vectors was eight units, in F14. Appendix Figure~\ref{fig:p2energy} locates these energies relative to the ground level.
\paperfigure{P2/plates/P2-F3_residual_margins.pdf}{Margins below the 196-unit decision boundary. Formal task selections span 84--176 above ground; the full formal return sets span 84--188 (shaded band). Five pilot tasks appear separately to the left.}{fig:p2margin}

Table~\ref{tab:p2service} counts each planned task once. All formal tasks yielded a strict solution, so the conditional regret column covers all twenty. Both endpoint definitions agree. Every formal energy lies above the certified ground level.
\begin{table}[!htb]
\centering
\caption{Attainment of the blind-pool proxy optimum by task or local restart block}\label{tab:p2service}
\begin{threeparttable}
\begin{tabular}{@{}lrrrrr@{}}
\toprule
Unit & Planned & Strict & Any hit & Selected hit & Regret$^{a}$\\
\midrule
Pilot task & 5 & 5 & 5 & 5 & 0\\
Formal task & 20 & 20 & 20 & 20 & 0\\
SA block & 5 & 5 & 5 & 5 & 0\\
\bottomrule
\end{tabular}
\begin{tablenotes}[flushleft]\footnotesize
\item[$a$] Maximum original regret of the energy-selected strict solution over units with a strict return. Each SA block has 100 restarts; internal service effort is undisclosed.
\end{tablenotes}
\end{threeparttable}
\end{table}
\FloatBarrier

\subsection{Model checks and sensitivity}\label{sec:p2checks}
\textbf{Where the decision was made.} Figure~\ref{fig:p2cloud} shows the decisive group directly: group 17 reaches its local minimum in every recorded return set, although the energy-selected formal vectors attain only 19--28 of the 40 group minima. The boundary case in Proposition~\ref{prop:p2margin} shows why a vote among groups could select the wrong material. Wei et al. describe the underlying class of optical coherent Ising machines~\cite{wei2026}; the measurements here concern the cloud optimization service.
\paperfigure{P2/plates/P2-F5_optimization_service.pdf}{Cloud optimization and group minima. Each heatmap cell counts local minima among ten returned solutions. The SPQC-550 is a generated illustration, not a photograph or identification of the executing backend. Service diagram adapted from the Kaiwu guide~\cite{kaiwu144}. The six auxiliary displays have no embedded task IDs; task associations come from separate records. Each uses its own vertical scale. The a.u. axis provides neither physical time nor evidence of evolution dynamics.}{fig:p2cloud}

\textbf{Sensitivity to the objective.} The predictor's error is measured against labeled materials; representation regret uses the frozen proxy table; integer residual uses the matrix energy. Appendix Figure~\ref{fig:p2cost} displays these quantities on separate scales. Across the 20 predictor tables, the adverse tail comprises candidates' largest regrets. In the blind pool both the mean-cost rule and the mean of the two largest member regrets select LaTiO$_2$N. Minimizing the single worst regret instead selects YTiO$_2$N. The full-pool comparison is in Appendix Figure~\ref{fig:p2tail}.

\textbf{Local controls.} At a budget of five restarts, 20 of 20 paired SA blocks found a strict hit with the encoded costs and six with the costs removed. Appendix Figure~\ref{fig:p2budget} gives the other budgets. Under a matched local schedule, one long chain yielded zero strict hits in 200 restarts and the 40-group encoding yielded 135 (Appendix Figure~\ref{fig:p2ablation}). These controls measure the local solver on this instance. The service data contain no matched hardware control for chain length or internal search effort.

\subsection{Direct answer to Problem 2}\label{sec:p2answer}
\begin{enumerate}
\item \textbf{Model.} The two-register QUBO in Eq.~\eqref{eq:p2qubo} selects an ordered A--B pair with the blind mask enforced. The cloud solver uses its decision-equivalent encoding of the blind dictionary in 40 domain-wall groups with one shared reference.
\item \textbf{Selection.} Choose La at A and Ti at B with $X=\mathrm{NO_2}$. LaTiO$_2$N uniquely minimizes the predicted violation, 1.839047, in both the 529-pair blind pool and the 2,704-pair full pool.
\item \textbf{Service result.} All 20 formal energy-selected strict solutions gave zero original regret. Their residuals, 84--176, lie below 196.
\end{enumerate}
